% !TEX root = ../notes.tex

\documentclass[../notes.tex]{subfiles}

\begin{document}

\section{March 4}
Today, we discuss a little about quantization.

\subsection{The Elliptic Calogero--Moser Hamiltonians}
Let's do another example of a twisted parabolic Hitchin system. Let $X$ be an elliptic curve with origin $0$, and let $P\subseteq\op{PGL}_n$ be the maximal parabolic subgroup of block upper-triangular matrices with sizes $(1,n-1)$. For the twisting, we are asking for a choice of scalar $c\in\CC$.

It turns out that a generic degree-$0$ rank $n$ vector bundle on $X$ is a direct sum of degree-$0$ line bundles. But line bundles on $X$ of degree zero are parameterized by points: for each $q\in X$, there is a line bundle
\[\mc L_q\coloneqq\OO(q)\otimes\OO(0)^{-1}.\]
This line bundle consists of those functions which look like $\theta(z-q)/\theta(z)$, where
\[\theta(z)\coloneqq\sum_{m,n}(-1)^ne^{2\pi inz+\pi n^2i\tau}\]
is the theta function on $X=\CC/(\ZZ+\tau\ZZ)$ (which is entire on $\CC$ and admitting only simple zeroes at $m+n\tau$, satisfying $\theta(z)=z+1$ and $\theta(z+\tau)=-e^{-z}\theta(z)$). In other words, we are looking at meromorphic functions with a single pole at $0$ and vanishes at $q$. Thus, a generic degree-$0$ rank $n$ vector bundle on $X$ looks like
\[\mc E=\bigoplus_{i=1}^n\mc L_{q_i}.\]
With $\mc E$ generic, we may even assume that the $q_\bullet$s are distinct. Now, $K_X$ is the trivial bundle (because $X$ has genus one), so a Higgs field $\mc E\to\mc E\otimes K_X$ is just the data of an endomorphism of $\mc E$. Such an endomorphism $\varphi$ corresponds to a matrix of endomorphisms $\varphi_{ij}\colon\mc L_{q_i}\to\mc L_{q_j}$, meaning that $\varphi_{ij}$ should be a global section of $\mc L_{q_j}\otimes\mc L_{q_i}^{-1}$. Keeping track of the parabolic structure, we are allowing $\varphi_{ij}$ to admit a simple pole at $0$ with some controlled residue. Accordingly, write
\[\varphi_{ij}=a_{ij}\cdot\frac{\theta'(0)\theta(z+q_j-q_i)}{\theta(z)\theta(q_j-q_i)}\]
for $i\ne j$, and $\varphi_{ii}=p_i$ is some momentum variable. Notably, $\varphi_{ij}$ has residue $a_{ij}$ at $0$, and $\varphi_{ii}$ has residue $0$ at $0$.

On the other hand, let's talk about what the parabolic structure gives us: a parabolic structure for our given parabolic amounts to the choice of a line in $\CC^n$. Now, for the twisting, $\op{Res}_0\varphi$ acts by $(n-1)c$ on this vector and on $-c$ on the quotient; by some normalization process using actions of the torus, we may even assume that this vector is $(1,\ldots,1)$. Explicitly, our twisting is given by the character $\chi$ on $\op{Lie}P$ given by
\[\begin{bmatrix}
	a & v \\
	0 & A
\end{bmatrix}\mapsto(n-1)ca-c\tr A.\]
Anyway, the point is that $\op{Res}_0\varphi+c\id$ should have rank $1$. But the $\varphi_{ii}$s have vanishing residue at $0$, so in the basis of the previous paragraph, we need the residue $a_{ij}$ to be $c$ everywhere. Thus,
\[\varphi_{ij}=c\cdot\frac{\theta'(0)\theta(z+q_j-q_i)}{\theta(z)\theta(q_j-q_i)}.\]
We are now ready to compute some polynomials in our Hitchin system. Doing a little calculation, we find that $H_2(z)=\frac12\tr\varphi^2$ is
\[\frac12\sum_ip_i^2+c^2\sum_{i\ne j}\frac{\theta'(0)^2\theta(z+q_i-q_j)\theta(z+q_j-q_i)}{\theta(z)^2}.\]
The latter sum turns out to be $\wp(z)-\wp(q_i-q_j)$. We are thus left with
\[c^2\cdot\frac{n(n-1)}2\wp(z)+\frac12\Bigg(\sum_ip_i^2-c^2\sum_{i\ne j}\wp(q_i-q_j)\Bigg).\]
The second sum is known as the elliptic Calogero--Moser Hamiltonian. The higher terms $\tr\varphi^m$ look like $\sum_ip_i^m$ plus some lower-order terms, which give the other coordinates. For example, one can calculate that $H_3$ looks like
\[\sum_ip_i^2+\sum_i(a_i(q)p_i+b(q)).\]

\subsection{Desiderata of Quantization}
We should start by explaining what our goal is when quantizing. A classical Hitchin system is an integrable system on $T^*Y$. In particular, we are trying to live in the ring $\OO(T^*Y)$, which is locally generated by $\OO(T)$ and some differential coordinates $p_i\coloneqq dy_i$. Quantization asks us to replace the variables $p_i$ by $\hbar\del_i$, where $\del_i\coloneqq\del/\del x_i$. Thus, $\OO(T^*Y)$ is quantized to quasiclassical differential operators $\mc D(Y)$, which is a Rees alegbra
\[\mc D(Y)=\sum_{j=0}^\infty\hbar^j\mc D_{\le j}(Y).\]
The point is that each derivative should have a formal factor of $\hbar$; sending $\hbar\del_i\to p_i$ and then send $\hbar\to0$ recovers the original.

To quantize our Hitchin system, we are thus interested in taking $H_\bullet$s from the original Hitchin system, and lifting them to $\widehat H_\bullet$s in $\mc D_{\hbar}(Y)$ which commute in the sense that $[\widehat H_i,\widehat H_j]=0$. Notably, we want $\widehat H_i$ to specialize to $H_i$ by the limiting process we just described.

Here is a basic example which tells us that such a thing might be possible.
\begin{example}[Gaudin system]
	For the Garnier system, we saw
	\[G_i=\sum_{j\ne i}\frac{(x_i-x_j)^2p_ip_j-2(x_i-x_j)(\lambda_ip_j-\lambda_jp_i)}{t_j-t_i}.\]
	Na\"ively, we may try to quantize this by setting
	\[\widehat G_i\coloneqq\sum_{j\ne i}\frac{\hbar^2(x_i-x_j)^2\del_i\del_j-(x_i-x_j)(\hbar\lambda_i\del_j-\hbar\lambda_j\del_i)}{t_i-t_j}.\]
	By some miracle, it turns out that these operators commute, which is not at all expected because $\del_i\del_j$ do not commute. In fact, the resulting system is independent of the ordering of the variables (up to some constants): one could even have used (for example) have put all $\del$s before the variables.
\end{example}
Let's explain where this construction comes from. Choose representations $V_1,\ldots,V_N$ of a simple Lie algebra $\mf g$. Set $\mc H\coloneqq V_1\otimes\cdots\otimes V_N$. There is a Casimir tensor $\Omega\in\left(S^2\mf g\right)^G$. One can then divide $\Omega$ into terms $\Omega_{ij}$ with $i\ne j$, meaning that the first component of the pure tensor $\Omega_{ij}$ acts on $V_i$, and the second acts on $V_j$. Then we set
\[\widehat G_i\coloneqq\sum_{j\ne i}\frac{\Omega_{ij}}{t_i-t_j},\]
which are operators on $\mc H$. It then turns out that having $[\Omega_{ij},\Omega_{k\ell}]=0$ and $[\Omega_{ij},\Omega_{ik}+\Omega_{jk}]=0$ for all $i$ and $j$ and $k$ implies that $[\widehat G_i,\widehat G_j]=0$ and descend to an action on $\mc H^G$.
\begin{example}
	With $\mf g=\mf{sl}_2$, take the representations $V_i$ to be $\CC[x]$, where $f$ acts by $-\del_x$, $e$ acts by $x^2\del_x-2\Lambda_ix$, and $h$ acts by $2x\del_x-\Lambda_i$.
\end{example}
\begin{remark}
	This process can only produce at most quadratic differential operators. In general, quantization must do something more complicated.
\end{remark}
Here is another example of quantization.
\begin{example}
	We can quantize the Calogero--Moser Hamiltonian
	\[H_2=\frac12\sum_ip_i^2-c^2\sum_{i\ne j}\wp(q_i-q_j)\]
	to
	\[\widehat H_2=\frac12\sum_i\hbar^2\del_i^2-c^2\sum_{i\ne j}\wp(q_i-q_j).\]
	After factoring out $\hbar^2$, one receives the ``quantum Calogero--Moser Hamiltonian'' $L_2$.
\end{example}
\begin{remark}
	It further turns out that there are operators $L_3,\ldots,L_n$ which commute and have $L_k$ is equal to $\sum_i\del_i^k$ up to lower-order terms. The commutativity of the classical system does not directly imply the commutativity of the higher-order terms, as one can see already from our choice of quantization of $H_3$.
\end{remark}

\subsection{Beilinson--Drinfeld Quantization}
Let's sketch the main ideas of how to quantize our classical Hitchin systems. For the time being, set $\hbar=1$. This requires us to redo much of the construction of the classical Hitchin system. Recall that the classical Hitchin system arose from Hamiltonian reduction: we wrote $\op{Bun}_G(X)$ as $G(R)\backslash G(K_x)/G(\OO_x)$, so Hamiltonian reduction granted us an identification
\[T^*\op{Bun}_G(X)=T^*G(K_x)/(G(R)\times G(\OO_x)).\]
Thus, we could find a Hitchin system by finding two-sided invariant functions on $T^*G(K_x)$ and then descend to the quotient.

Thus, we must replace $T^*\op{Bun}_G(X)$ by $\mc D_{\hbar}(\op{Bun}_G(X))$, and then we need to take $2$-sided invariant elements in $\mc D_{\hbar}(G(K_x))$ and descend to the quotient. This will not quite work, but it is a reasonable first attempt.
\begin{example}
	To get our bearings, we work in the case where $X$ is a Lie group $K$. Then left-invariant differential operators in $\mc D(K)$ are given by $U\mf k$, where $\mf k\coloneqq\op{Lie}K$, so the two-sided invariant differential operators in $\mc D(K)$ are given be the center $Z\mf k$. Now, our $H_2=\frac12\tr\varphi^2$ looks like $\varphi=\sum_n\varphi z^{-n-1}\,dz$, where $\varphi_n\in\mf g$. Now, letting $\{a_i\}$ be an orthonormal basis of $\mf g$ for the Killing form so that $\varphi=\sum_i\varphi^ia_i$, we find that
	\[\frac12\tr\varphi^2=\sum_nT_nz^{-n-2}\,(dz)^2,\]
	where $T_n=\sum_{m,i}\varphi^i_m\varphi^{i}_{n-m}$. (To make sense of this, we can view the $T_n$s as some quadratic elements of $U\mf g((t))$.) For commutativity, we require $[\varphi^j_\ell,T_n]=0$; classically, the Poisson bracket does vanish, but we no longer achieve commutativity because there are ordering issues. After a calculation, one finds $[\varphi^j_\ell,T_i]=\ell h^\lor\varphi^j_{\ell+n}$, where $h^\lor$ is the dual Coxeter number.
\end{example}
There is no obvious fix for this, and indeed, there shouldn't be: Beilinson and Drinfeld showed that $\op{Bun}_G(X)$ admits no global nonconstant differential operators. The solution is to consider differential operators on $\frac12$-densities.
\begin{remark}
	In fact, this is more natural from the perspective of quantum mechanics: the quantization of $T^*M$ should be a Hilbert space. But to have a canonical pairing, we need to use half-densities instead of functions.
\end{remark}
Accordingly, we should replace $\mc D(G(K))$ should be replaced by differential operators on some line bundle. (Here, $G(K)$ is the loop group.) Luckily, there is a suggestion of such a line bundle: the Kac--Moody group $\widehat G$ fits into a central extension
\[1\to\CC^\times\to\widehat G\to G((t))\to1,\]
and the Lie algebra of $\widehat G$ is some Kac--Moody algebra. It turns out that the canonical bundle of $\op{Bun}_G(X)$ is $\mc L^{-2h^\lor}$, where $\mc L$ is the associated line bundle on the loop group (which notably descends to $\op{Bun}_G$). The moral is that we can replace $U\mf g((t))$ with the enveloping algebra of this Kac--Moody algebra above up to identifying $K$ with $-h^\lor$. Doing this quotient expands the center significantly (called the Feigen--Frenkel center), so we get our integrable system by redoing the whole system.

\end{document}